% BANKS-SOURCE: pdf=49 print=39 kind=section id=4.2
\subsection{Massless particles with helicity}

% BANKS-SOURCE: pdf=49 print=39 kind=prose
The helicity discussion in this subsection is the four-dimensional specialization. In
$D>4$, the massless little group is $\mathrm{ISO}(D-2)$, so a single helicity label does
not describe all of its finite-dimensional representations.

When we try to generalize these considerations to massless particles, we run into a
surprise. We can always transform a null momentum to a frame where
$p^\mu=(E,\mathbf p_\star)$ with $\mathbf p_\star=(0,0,E)$.
This is invariant under rotations around the 3-direction. The combination of transverse
boost and rotation generators $T_i=K_i+\epsilon_{ij}J^j$ also leaves the null momentum invariant.
Together with the transverse rotation these form the two-dimensional Euclidean group.
This group has infinite-dimensional ``continuous spin'' representations unless $T_i=0$.
There are no such infinitely degenerate massless particles in nature, and one cannot
make a local field theory which includes them.\footnote[1]{T. Banks (1973), \emph{Continuous Spin Representations and Locality} (unpublished).}
The remaining representations are characterized by the eigenvalue of $J_3$, called the helicity, which must be quantized in
half-integer units in order to have a local field theory.\footnote[2]{You will prove the assertions on this page in the problems at the end of the chapter.}
Furthermore, local field theory requires that a helicity value $h$ be accompanied by $-h$. Indeed, if this were not the
case we could trace the causal order of emission and absorption processes between
space-like separated points by following the helicity flow. However, the causal order
can be changed by doing boosts along the momentum of the massless particle, which
preserve helicity. On a technical level we find that we need both helicities (and helicity
quantization) in order to construct a local field.

% BANKS-SOURCE: pdf=49 print=39 kind=display id=4.2a
A general null momentum $p^\mu=(|\mathbf p|,\mathbf p)$ can be obtained from the canonical one
$p_\star{}^\mu=(E,\mathbf p_\star)$, where $\mathbf p_\star=(0,0,E)$, by performing a boost
of rapidity $\ln(|\mathbf p|/E)$ along the 3-direction, followed by a rotation
of the 3-direction into the direction $\mathbf p/|\mathbf p|$. Call the product of these two transformations
$L(p)$. By analogy with the massive case, we can derive the Lorentz transformation rule
for massless single-particle states. Here $\omega_{\mathbf p}=|\mathbf p|$, and the rule is

% BANKS-SOURCE: pdf=49 print=39 kind=equation id=4.2
\begin{equation}
  U(\Lambda)\ket{p,h}=\sqrt{\frac{\omega_{\boldsymbol{\Lambda p}}}{\omega_{\mathbf p}}}
  U\bigl(L^{-1}(\Lambda p)\Lambda L(p)\bigr)\ket{p,h}
  =\ee^{\ii h\Theta(\Lambda,p)}\ket{\Lambda p,h}.
  \label{eq:4.2}
\end{equation}
The second equality follows because $L^{-1}(\Lambda p)\Lambda L(p)$ is in the stability subgroup of
$(|\mathbf p|,\mathbf p)$ and so acts on the states like a rotation around the direction of motion. $\Theta(\Lambda,p)$
is the angle of that rotation.

% BANKS-SOURCE: pdf=49 print=39 kind=prose
In Problem 4.3, you will show that the construction of local fields from the creation
and annihilation operators of massless particles requires that $h$ be quantized in half-
integer multiples and that a particle of helicity $-h$ exists for every particle of helicity $h$.
For particles that participate in parity-conserving interactions, it is conventional to call
the state with helicity $-h$ another spin state of the same particle, whereas for neutrinos,
which have only parity-violating interactions, it is called the anti-neutrino.

% BANKS-SOURCE: pdf=49 print=39 kind=prose
Note that for $|h|>\frac12$ there is a discontinuity in the number of degrees of freedom
of a massive particle of spin $h$ and a massless particle with helicity $h$. The ``missing''
degrees of freedom are called \emph{longitudinal}. A field theory of massive particles of spin
$h$, which has a smooth massless limit, must acquire some sort of symmetry principle in
the massless limit, which decouples the unwanted longitudinal states. By dimensional

% BANKS-SOURCE: pdf=50 print=40 kind=prose
analysis, going to large momentum must be related to the zero-mass limit, so a theory
with good high-energy behavior must have a similar decoupling principle. We will see
that this requirement is the origin of gauge invariance, both for Maxwell’s field and for
much of the rest of the \emph{standard model} of particle physics. To that end, we proceed to
the field theory of massive spin-1 particles, after which we will take the massless limit.
